\section{System overview}
\label{sec:system}

Figure~\ref{fig:architecture} shows AutoTerm-SST's online terminology-selection
loop around RASST, whose models stay unchanged. RASST translates 1.92-second
audio chunks with a speech encoder, a cross-modal term retriever, and the
Qwen3-Omni LLM. The encoder maps chunk $t$ to features $h_t$, the retriever
builds a term map $G_t$, and the LLM translates $h_t$ with $G_t$ in its prompt.
Without a source transcript from RASST, the retriever matches speech directly
to glossary source terms by maximum similarity, or MaxSim, between overlapping
speech windows and pre-encoded term vectors. Terms scoring above 0.78 become
candidates. The retriever adds up to ten bilingual term pairs to $G_t$.
AutoTerm-SST changes only retrieval scope. Its router uses each translation
to select up to four topic slices for later searches.

Keeping adaptation outside the translation model makes the system pluggable.
One streaming protocol supports a retrieval-free translator, fixed-glossary
RASST, and AutoTerm-SST. Each partial translation sends the retrieved pairs,
their source slices, and the working set to the interface.
